\section*{Cel i~zakres pracy}
\addcontentsline{toc}{section}{Cel i~zakres pracy}
\label{chap:cel}

\subsection*{Cel główny pracy}

Celem pracy jest zaprojektowanie i~zaimplementowanie prototypu aplikacji webowej RiskScoop AI -- systemu analizy ryzyka powodziowego dla infrastruktury krytycznej, wykorzystującego agenta AI do przetwarzania zapytań w~języku naturalnym i~generowania analiz geoprzestrzennych. System ma zastosowanie w~obszarze Rapid Mapping, łącząc interfejs języka naturalnego z~oceną ryzyka powodziowego.

\subsection*{Cele szczegółowe}

W~ramach realizacji celu głównego wyznaczono następujące cele szczegółowe:

\begin{enumerate}
    \item \textbf{Analiza literatury} obejmująca:
    \begin{itemize}
        \item metody uczenia maszynowego w~ocenie ryzyka powodziowego,
        \item agentów AI wykorzystujących interfejs języka naturalnego w~analizie geoprzestrzennej (ChatGeoAI, LLM-Find, GeoGPT),
        \item systemy Rapid Mapping oraz Copernicus GloFAS,
        \item otwarte dane geoprzestrzenne (Overture Maps Foundation),
        \item technologie front-end dla aplikacji webowych (Mapbox GL JS).
    \end{itemize}

    \item \textbf{Przeprowadzenie analizy konkurencyjnych rozwiązań} na rynku systemów analizy geoprzestrzennej z~interfejsem języka naturalnego oraz platform specjalizujących się w~ocenie ryzyka powodziowego, w~celu identyfikacji luk rynkowych i~możliwości dla RiskScoop AI.

    \item \textbf{Zaprojektowanie architektury systemu} RiskScoop AI obejmującej:
    \begin{itemize}
        \item komponent agenta AI do przetwarzania języka naturalnego,
        \item moduł integracji z~bazą Overture Maps,
        \item moduł integracji z~systemem Copernicus GloFAS,
        \item moduł kwantyfikacji ryzyka powodziowego,
        \item interfejs użytkownika z~interaktywną wizualizacją mapową.
    \end{itemize}

    \item \textbf{Zaimplementowanie prototypu systemu} wykorzystującego:
    \begin{itemize}
        \item Vanilla JavaScript i~Mapbox GL JS do wizualizacji geoprzestrzennej,
        \item backend obsługujący przetwarzanie zapytań i~operacje geoprzestrzenne,
        \item integrację z~modelami języka naturalnego (LLM).
    \end{itemize}

    \item \textbf{Przeprowadzenie testów funkcjonalnych i~wydajnościowych} zaimplementowanego prototypu, weryfikujących:
    \begin{itemize}
        \item poprawność interpretacji zapytań w~języku naturalnym,
        \item dokładność identyfikacji obiektów infrastrukturalnych,
        \item poprawność kwantyfikacji ryzyka powodziowego,
        \item wydajność systemu przy różnych obciążeniach.
    \end{itemize}
\end{enumerate}
